\section{Results on Fresh Shared Routes}
\begin{table}[t]
\caption{Strict joint success on the new shared routes (percent).}
\label{tab:jointresults}\centering\small
\begin{tabular}{lrrr}\toprule
Controller & Standard & Longer & Per split\\\midrule
PPO original (V1.5) & 95.7 & 85.8 & 1,000\\
PPO full (V2) & 95.7 & 87.3 & 1,000\\
PPO arrival (V2) & 94.9 & 85.1 & 1,000\\
Goal RSS & 94.0 & 84.5 & 200\\
Goal radio & 96.5 & 88.5 & 200\\
Joint one step & 96.0 & 90.5 & 200\\
Joint three steps & 92.0 & 90.5 & 200\\
\bottomrule\end{tabular}
\end{table}

\begin{figure*}[t]
\centering\includegraphics[width=.97\textwidth]{reward_comparison_success.pdf}
\caption{Joint success on fresh shared routes. Bars aggregate episodes; dots
show individual learned training seeds, not confidence intervals. Deterministic
controllers have one evaluation per route.}
\label{fig:success}
\end{figure*}

\subsection{Primary Reward Contrast}
Full achieves \FullStd\% standard success versus \OriginalStd\% for V1.5.
The paired difference is \PrimaryDiff\ percentage points, 95\% interval
\PrimaryCI. \PrimaryInterpretation{} On longer routes the difference is
\LongDiff\ points, interval \LongCI. \LongInterpretation{}
Every training seed appears in Table~\ref{tab:seeds}.

\begin{table}[t]
\caption{First failures: RSS / buffer / timeout counts. No energy or boundary failures occurred.}
\label{tab:failures}\centering\small
\begin{tabular}{lrr}\toprule
Controller & Standard & Longer\\\midrule
PPO original (V1.5) & 40 / 3 / 0 & 136 / 6 / 0\\
PPO full (V2) & 34 / 9 / 0 & 112 / 15 / 0\\
PPO arrival (V2) & 51 / 0 / 0 & 147 / 0 / 2\\
Goal RSS & 12 / 0 / 0 & 31 / 0 / 0\\
Goal radio & 7 / 0 / 0 & 21 / 2 / 0\\
Joint one step & 8 / 0 / 0 & 18 / 0 / 1\\
Joint three steps & 4 / 0 / 12 & 11 / 0 / 8\\
\bottomrule\end{tabular}
\end{table}

\subsection{Failures and Communication Tradeoffs}
Table~\ref{tab:failures} reports first failure reasons. Episodes stop on the
first violation; hypothetical later failures are unknown. Every successful
flight has zero sampled RSS outage and zero dropped bits.

\begin{table*}[t]
\caption{Full minus original on 925 standard and 781 longer common successful pairs. Secondary intervals are descriptive and unadjusted.}
\label{tab:paired}\centering\small
\begin{tabular}{llrrrr}\toprule
Split & Metric & Original & Full & Difference & 95\% interval\\\midrule
Standard & Flight time (s) & 31.252 & 28.112 & -3.139 & $[-3.509, -2.808]$\\
Standard & Accumulated radio cost & 7.094 & 5.309 & -1.785 & $[-2.194, -1.389]$\\
Standard & Delay proxy (s) & 1.942 & 5.827 & +3.886 & $[2.731, 5.066]$\\
Standard & Handovers & 6.843 & 0.443 & -6.400 & $[-8.501, -4.366]$\\
Standard & Interference ($\mu$W) & 0.753 & 0.748 & -0.004 & $[-0.017, 0.009]$\\
Standard & Energy proxy (kJ) & 7.651 & 7.788 & +0.137 & $[0.115, 0.157]$\\
Standard & SINR (dB) & -8.322 & -9.305 & -0.983 & $[-1.483, -0.488]$\\
Standard & Path length (m) & 547.596 & 551.204 & +3.608 & $[2.947, 4.273]$\\
Standard & Network overrides & 0.107 & 0.478 & +0.371 & $[0.243, 0.504]$\\
Standard & Motion overrides & 0.008 & 0.016 & +0.009 & $[-0.004, 0.028]$\\
Longer & Flight time (s) & 65.394 & 60.145 & -5.250 & $[-5.815, -4.804]$\\
Longer & Accumulated radio cost & 16.800 & 13.803 & -2.997 & $[-3.765, -2.257]$\\
Longer & Delay proxy (s) & 2.857 & 6.818 & +3.961 & $[2.657, 5.145]$\\
Longer & Handovers & 15.826 & 1.528 & -14.298 & $[-19.543, -9.101]$\\
Longer & Interference ($\mu$W) & 0.794 & 0.791 & -0.003 & $[-0.015, 0.009]$\\
Longer & Energy proxy (kJ) & 18.523 & 18.957 & +0.434 & $[0.394, 0.476]$\\
Longer & SINR (dB) & -8.773 & -9.957 & -1.184 & $[-1.652, -0.712]$\\
Longer & Path length (m) & 1345.958 & 1349.568 & +3.610 & $[2.896, 4.321]$\\
Longer & Network overrides & 0.467 & 1.757 & +1.289 & $[0.882, 1.637]$\\
Longer & Motion overrides & 0.024 & 0.044 & +0.019 & $[-0.010, 0.062]$\\
\bottomrule\end{tabular}
\end{table*}

Full versus original secondary comparisons contain \PairStd\ common successful
standard seed/route pairs and \PairLong\ longer pairs.
Table~\ref{tab:paired} reports conditional paired values and intervals.
\RadioInterpretation{} A favorable aggregate cost is not interchangeable with
delay, interference, handovers or completion. Different successful subsets
must not be treated as matched counterfactual outcomes.

\subsection{Other Controls}
Full versus arrival completion differences are \ArrivalDiffStd\ points
\ArrivalCIStd\ on standard routes and \ArrivalDiffLong\ points
\ArrivalCILong\ on longer routes. Both arms enforce strict constraints.
Goal radio attains \RadioStd\% and \RadioLong\% joint success.
\ReferenceInterpretation{} These references expose performance available from
the shared model rather than demonstrating learned superiority.

\begin{figure*}[t]
\centering\includegraphics[width=.97\textwidth]{reward_comparison_example.pdf}
\caption{First declared standard test route, seed 2101 for learned policies.
Original reward, full reward and Goal radio share initial conditions. Selection
is by index; aggregate results determine the conclusions.}
\label{fig:example}
\end{figure*}

\section{Interpretation and Limitations}
The contrast isolates the reward package conditional on V2. It does not identify
individual terms or behavior without the navigation prior and prospective
filter. These components may compensate for reward deficiencies. Removing them
would require a separate experiment and fresh evaluation.

Original speed limits, motion, discrete action parameterization, continued
training after arrival and unknown implementation details differ. Reconstructing
the reward on V2 cannot identify the cause of the thesis's wandering
\cite[p.~21, Sec.~8]{thesis}. Exact dataset version and sentinel meaning remain
unverified. One map, static load, ideal service and uncalibrated energy limit
inference. No physical continuity, packet latency, collision safety or deployment
claim follows. A successful arrival can still retain queued data. Every table
here uses the newly declared route seeds; earlier V2 results remain separate.

\section{Conclusion}
\ConclusionResult{} The experiment compares the written original reward and a
replacement under identical sampled constraints and control machinery, retaining
failed missions and communication tradeoffs. Further claims require independent
maps, traffic and controller ablations rather than tuning on these test routes.

\clearpage
\appendices
\section{Artifacts and Reproduction}
The repository includes all 15 final checkpoints, source, configuration, raw
records, analysis, source snapshots and this paper:
\url{https://github.com/unworthyzeus/UAV_joint_handover_trajectory_optimization}.
Ten are unchanged V2 policies and five are new V1.5 policies. The private
dataset belongs at \texttt{dataset/Barcelona\_dataset\_January.h5}; the setup
guide records its checksum, dependency versions, exact commands and outputs.
Note 36 is the frozen protocol, note 37 records results, and note 38 records
delivery validation. The README documents the exact V2 differences from the
TFM with page locators. This is a research draft, not an accepted publication.

\section{Every Seed and Observation Layout}
\par\noindent\begin{minipage}{\columnwidth}
\captionof{table}{Every learned seed: joint success percentages.}
\label{tab:seeds}\centering\small
\begin{tabular}{lrrr}\toprule
Arm & Seed & Standard & Longer\\\midrule
PPO original (V1.5) & 2101 & 96.5 & 85.5\\
PPO original (V1.5) & 2102 & 98.0 & 89.0\\
PPO original (V1.5) & 2103 & 95.0 & 85.0\\
PPO original (V1.5) & 2104 & 94.0 & 85.5\\
PPO original (V1.5) & 2105 & 95.0 & 84.0\\
PPO full (V2) & 2101 & 97.0 & 86.5\\
PPO full (V2) & 2102 & 95.0 & 87.5\\
PPO full (V2) & 2103 & 95.5 & 85.5\\
PPO full (V2) & 2104 & 95.5 & 87.5\\
PPO full (V2) & 2105 & 95.5 & 89.5\\
PPO arrival (V2) & 2101 & 95.5 & 83.5\\
PPO arrival (V2) & 2102 & 96.5 & 84.5\\
PPO arrival (V2) & 2103 & 94.0 & 86.5\\
PPO arrival (V2) & 2104 & 93.5 & 82.5\\
PPO arrival (V2) & 2105 & 95.0 & 88.5\\
\bottomrule\end{tabular}
\end{minipage}\par

\par\noindent\begin{minipage}{\columnwidth}
\captionof{table}{Identical V2 observation layout in all three learned arms.}
\centering\small
\begin{tabular}{ll}\toprule
Indices & Meaning\\\midrule
0--1 & Position / map extent\\
2--3 & Unit goal direction\\
4 & Goal distance / 1,000 m\\
5--6 & Radial/lateral velocity / 25 m/s\\
7--8 & Speed / 25; elapsed steps / 200\\
9--10 & Queue and consumed energy fractions\\
11--13 & Load phase, RBG and serving SINR\\
14--18 & Candidate RSS values\\
19--28 & Candidate relative coordinates\\
29--33 & Candidate station indices\\
34--94 & Log capacity for 61 options\\
95--155 & Network mask\\\bottomrule
\end{tabular}
\end{minipage}\par
